Instance normalization (RevIN and its successors) has become the default level-handling layer of multi-step forecasters, and every member of the family estimates the level it restores \emph{endogenously}, from the target series' own past. We show that this default has an applicability boundary, locate it, and give a coordinate that separates its two sides. A closed-form accounting in a linear-Gaussian instance model makes the trade explicit: instance normalization is immune to level variance and to covariate-orthogonal distribution shift---the property its designers intended---but pays the full within-horizon level displacement plus a window-noise floor, so its advantage falls as the exogenous predictability of the level rises and as the forecast horizon grows. Empirically, on a pre-registered grid of 1{,}794 runs---eight datasets (four exogenously driven energy series, four standard benchmarks), four backbones, three horizons, five seeds---the predicted sign of the effect was correct on all eight datasets, and the record is unchanged under MAE and under MASE. Two results separate the \emph{layer} from the information it is denied. Under information parity, with identical past and future covariates supplied to every arm, instance normalization is worse in mean MSE than plain global $z$-scoring on all five covariate-capable backbones---a function-class constraint, not an information advantage. And in the main table, three different covariate-conditioned level models---a first-stage regression used alone, a classical dynamic regression, and conditional normalization---beat the best instance-normalized arm in 11 of 12 comparisons on the exogenously driven group and in 0 of 12 on the standard benchmarks (10 of 12 and 0 of 12 under MASE), against seasonal-naive and climatology anchors. The result is therefore about the layer, not about any one replacement for it: when the level is set by observable drivers, subtracting the window mean discards recoverable signal, and any means of conditioning the level on covariates repairs it. Variance attribution locates the effect where a boundary should sit---the normalization$\times$dataset interaction explains $47.1\%$ of log-MSE variance against $0.6\%$ for all backbone-related terms. We report the out-of-sample $R^2$ of window-mean levels on covariates (the Level Predictability Score, $\lps$) as a pre-training coordinate that separates the two regimes, and we are explicit about what it does not do: our eight datasets are bimodal in $\lps$, the threshold is stress-tested in effectively one case, and within the exogenous regime $\lps$ does not order the magnitude of the effect. $\lps$ is demonstrated as a regime classifier, not as a calibrated dial.
